% ============================================================
% git-init-existing-folder.tex — A SUBDOCUMENT
% Compiles standalone ("pdflatex git-init-existing-folder.tex")
% using main.tex's preamble, OR gets inserted into main.tex via
% \subfile{git-init-existing-folder}.
% ============================================================

\documentclass[../main.tex]{subfiles}

\begin{document}

\section{Creating a New GitHub Repository from an Existing Folder}

\subsection{Overview}
This guide covers turning a folder you already have on your
computer (with existing files) into a Git repository, then
pushing it to a brand-new repository on GitHub.

\subsection{Step 1: Open Git Bash in the Folder}
\begin{enumerate}[leftmargin=*]
    \item Navigate to the folder in Windows File Explorer.
    \item Right-click inside the folder and select
    \textbf{Git Bash Here}.
\end{enumerate}

Alternatively, open Git Bash and use \texttt{cd} to navigate there:

\begin{lstlisting}
cd /c/Users/yourname/Documents/your-project
\end{lstlisting}

\subsection{Step 2: Initialize the Folder as a Git Repository}
\begin{lstlisting}
git init
\end{lstlisting}

This creates a hidden \texttt{.git} folder inside your project,
which is what turns an ordinary folder into a Git repository.

\subsection{Step 3: Check the Status}
\begin{lstlisting}
git status
\end{lstlisting}

This lists all the files in the folder as ``untracked'' — Git
sees them but isn't tracking changes to them yet.

\subsection{Step 4: Stage and Commit Your Files}
\begin{lstlisting}
git add .
\end{lstlisting}

The \texttt{.} stages every file in the folder (and subfolders)
for the next commit.

\begin{lstlisting}
git commit -m "Initial commit"
\end{lstlisting}

The \texttt{-m} flag lets you attach a short message describing
the commit, in quotes.

\subsection{Step 5: Set Your Default Branch Name (Optional)}
Modern GitHub repositories default to a branch named
\texttt{main}. If your local repo defaulted to \texttt{master}
instead, rename it:

\begin{lstlisting}
git branch -M main
\end{lstlisting}

\subsection{Step 6: Create a New, Empty Repository on GitHub}
\begin{enumerate}[leftmargin=*]
    \item Log in to GitHub and click the \textbf{+} icon in the
    top-right corner, then select \textbf{New repository}.
    \item Enter a \textbf{Repository name}.
    \item Choose \textbf{Public} or \textbf{Private}.
    \item \textbf{Important:} Do \emph{not} check any of the boxes
    to add a README, \texttt{.gitignore}, or license. Since you
    already have files locally, starting the GitHub repo empty
    avoids merge conflicts.
    \item Click \textbf{Create repository}.
\end{enumerate}

\subsection{Step 7: Link Your Local Folder to the GitHub Repository}
After creating the repository, GitHub displays a remote URL.
Copy the \textbf{SSH} URL (it should look like
\texttt{git@github.com:username/repository.git}) and run:

\begin{lstlisting}
git remote add origin git@github.com:username/repository.git
\end{lstlisting}

Replace \texttt{username/repository} with your actual GitHub
username and repository name.

\subsection{Step 8: Push Your Files to GitHub}
\begin{lstlisting}
git push -u origin main (or whatever you called it)
\end{lstlisting}

The \texttt{-u} flag sets \texttt{origin main} as the default
remote and branch, so future pushes only require:

\begin{lstlisting}
git push
\end{lstlisting}

\subsection{Step 9: Verify}
Refresh the repository page on GitHub. Your files should now
appear there, matching what's in your local folder.

\subsection{Note on SSH}
This guide assumes SSH is already set up between your computer
and GitHub. If \texttt{git push} asks for a username and password
instead of working automatically, you'll need to add an SSH key
to your GitHub account first.

\end{document}